% from ega0-17-fr.tex
\oldpage[0\textsubscript{IV-1}]{46}
de (17.2.13.2) est égal à
$\operatorname{Tor}^B_i(M\otimes_A B,N\otimes_A B)$~; l'hypothèse sur $B$
entraîne que ce $B$-module est nul pour $i>n$ (17.2.2, (ii)), d'où la
conclusion.
\end{proof}

\begin{env}[17.2.14]
\label{0.17.2.14-fr}
Soient $(X,\mathscr{O}_X)$ un espace annelé, $\mathscr{F}$ un
$\mathscr{O}_X$-Module~; on appelle \emph{dimension projective} (resp.
\emph{injective}) \emph{ponctuelle} de $\mathscr{F}$ et on note
$\operatorname{dim.proj}(\mathscr{F})$ (resp.
$\operatorname{dim.inj}(\mathscr{F})$) le nombre
$\sup_{x\in X}(\operatorname{dim.proj}(\mathscr{F}_x))$ (resp.
$\sup_{x\in X}(\operatorname{dim.inj}(\mathscr{F}_x))$. On appelle
\emph{dimension cohomologique ponctuelle de $X$} et on note
$\operatorname{dim.coh}(X)$ le nombre
$\sup_{x\in X}(\operatorname{dim.coh}(\mathscr{O}_x))$. Il résulte de
(17.2.5) et (17.2.10) que lorsque $A$ est un anneau noethérien, $M$ un
$A$-module de type fini et $X=\operatorname{Spec}(A)$, on a
$\operatorname{dim.proj}(\widetilde{M})=\operatorname{dim.proj}(M)$,
$\operatorname{dim.inj}(\widetilde{M})=\operatorname{dim.inj}(M)$ et
$\operatorname{dim.coh}(X)=\operatorname{dim.coh}(A)$. On appelle
\emph{dimension projective} (resp. \emph{injective}) \emph{de $\mathscr{F}$ en
un point $x\in X$} la dimension projective (resp. injective) de
$\mathscr{F}_x$, \emph{dimension cohomologique de $X$ en un point $x$} la
dimension cohomologique de $\mathscr{O}_x$.
\end{env}

\begin{proposition}[17.2.15]
\label{0.17.2.15-fr}
Soient $X$, $Y$ deux espaces annelés en anneaux locaux noethériens,
$f:X\longrightarrow Y$ un morphisme plat. Si
$\operatorname{dim.coh}(X)\leq n$, alors $Y$ est de dimension cohomologique
$\leq n$ en tout point de $f(X)$.
\end{proposition}

\begin{proof}
Cela résulte aussitôt de (17.2.13).
\end{proof}

\subsection{Théorie cohomologique des anneaux réguliers}
\label{subsection:0.17.3-fr}

\begin{theorem}[17.3.1]
\label{0.17.3.1-fr}
\textnormal{(Hilbert-Serre)} Soit $A$ un anneau local noethérien. Pour que $A$
soit de dimension cohomologique finie, il faut et il suffit que $A$ soit
régulier~; on a alors
\begin{equation}
\label{0.17.3.1.1-fr}
\tag{17.3.1.1}
  \operatorname{dim.coh}(A)=\dim(A).
\end{equation}
\end{theorem}

\begin{proof}
Supposons $A$ régulier~; soient $\mathfrak{m}$ son idéal maximal,
$\mathbf{x}=(x_i)_{1\leq i\leq n}$ un système régulier de paramètres de $A$
(17.1.6)~; considérons le complexe de l'algèbre extérieure
$K_\bullet(\mathbf{x})$ (\textbf{III}, 1.1.1), qui est formé de $A$-modules
libres, avec $K_i(\mathbf{x})=0$ pour $i>n$~; comme $(x_i)$ est une suite
régulière, on a $H_i(K_\bullet(\mathbf{x}))=0$ pour $i>0$ (\textbf{III},
1.1.4 et 1.1.3.3) et
$H_0(K_\bullet(\mathbf{x}))=A/(x_1A+\cdots+x_nA)=A/\mathfrak{m}=k$~; les
$K_i(\mathbf{x})$ forment par suite une \emph{résolution libre} de $k$ de
longueur $n$. Or, le fait que les $x_i$ appartiennent à $\mathfrak{m}$
entraîne aussitôt que dans le complexe
$K_\bullet(\mathbf{x},k)=K_\bullet(\mathbf{x})\otimes_A k$, l'opérateur bord
est nul en toutes dimensions, si bien que l'on a, par définition,
$\operatorname{Tor}^A_i(k,k)=\bigwedge^i(k^n)$~; l'égalité (17.3.1.1) résulte
donc aussitôt de (17.2.11) (ce résultat est essentiellement le «~théorème des
syzygies~» de Hilbert).

Montrons maintenant que si $A$ est un anneau local noethérien,
$\mathfrak{m}$ son idéal maximal, et si
$\operatorname{dim.proj}_A(\mathfrak{m})$ est \emph{finie}, alors $A$ est
régulier, ce qui achèvera de démontrer (17.3.1). Nous procéderons par
récurrence sur $n=\operatorname{rg}_k(\mathfrak{m}/\mathfrak{m}^2)$. Pour
$n=0$, on a $\mathfrak{m}=0$ et l'assertion est triviale.
\end{proof}

\begin{lemma}[17.3.1.2]
\label{0.17.3.1.2-fr}
\textnormal{(Nagata)} Si tout élément de
$\mathfrak{m}-\mathfrak{m}^2$ est diviseur de zéro dans $A$, il existe
$c\neq0$ dans $A$ tel que $c\mathfrak{m}=0$ (autrement dit on a
$\mathfrak{m}\in\operatorname{Ass}(A)$).
\end{lemma}

\begin{proof}
On peut se borner au cas où $\mathfrak{m}\neq0$, donc
$\mathfrak{m}\neq\mathfrak{m}^2$. L'hypothèse entraîne que
$\mathfrak{m}-\mathfrak{m}^2$ est contenu dans la réunion des idéaux
$\mathfrak{p}_i$ de $\operatorname{Ass}(A)$ (Bourbaki,
\emph{Alg.~comm.}, chap.~IV, \S~1, no~1, cor.~3 de la prop.~2)~; donc
$\mathfrak{m}$ est contenu dans la réunion de $\mathfrak{m}^2$
